\section{Dataset}
\label{sec:dataset}

\subsection{MovieLens 100K Overview}

The MovieLens 100K dataset is a widely-used benchmark in recommendation systems research. It was collected 
by the GroupLens research group at the University of Minnesota over a seven-month period (Sept 19th, 1997 -- 
April 22nd, 1998). The dataset consists of approximately 100,000 ratings provided by around 1,000 users on 
1,700 movies.

\textbf{Data Distribution:}
\begin{itemize}
  \item \textbf{Number of users}: 943
  \item \textbf{Number of items (movies)}: 1,682
  \item \textbf{Number of ratings}: 100,000
  \item \textbf{Rating scale}: 1--5 (integer values)
  \item \textbf{Sparsity}: $\frac{100,000}{943 \times 1,682} \approx 6.3\%$ (about 93.7\% of the matrix is missing)
\end{itemize}

The high sparsity is typical of real-world recommendation problems and is central to the challenge: 
predicting ratings for user-item pairs that have never been observed.

\subsection{Data Format and Preprocessing}

The MovieLens 100K dataset is provided in a tabular format. Each row represents a single rating observation, 
with fields:
\begin{itemize}
  \item User ID: integer identifier (1--943)
  \item Item ID: integer identifier (1--1,682)
  \item Rating: integer in $\{1, 2, 3, 4, 5\}$
  \item Timestamp: Unix timestamp of the rating event
\end{itemize}

The \code{Surprise} library handles dataset loading and preprocessing automatically. Specifically:

\begin{enumerate}
  \item \textbf{String conversion}: User and item IDs are converted to strings for API compatibility.
  \item \textbf{Train-test splits}: The dataset provides standard cross-validation folds (u1--u5) and 
  alternative train-test splits (ua, ub) for benchmarking.
  \item \textbf{No additional preprocessing}: Ratings are used as-is. No normalization, outlier removal, 
  or feature engineering is applied in the base implementation.
\end{enumerate}

\subsection{Statistical Characteristics}

Analysis of the MovieLens 100K dataset reveals:

\textbf{Rating Distribution:}
The ratings follow a roughly normal distribution with slight skew toward positive ratings. The mean rating 
is approximately 3.5, indicating users tend to rate items they watch (selection bias).

\textbf{User Activity:}
Users vary dramatically in activity level. Some users have provided hundreds of ratings, while others 
contributed only a handful. This heterogeneity impacts model training: high-activity users provide more 
signal, while low-activity (cold-start) users are harder to model.

\textbf{Item Popularity:}
Similarly, movies vary in the number of ratings they receive. Popular movies have many ratings, enabling 
reliable factor estimation, while niche movies have few ratings and are prone to cold-start effects.

\textbf{Temporal Characteristics:}
Though the dataset is small by modern standards and temporally localized to 1997--1998, the timestamps allow 
for temporal analysis. For this project, temporal information is not utilized, though it could support 
cold-start improvements or drift detection in production.

\subsection{Data Assumptions and Limitations}

Important assumptions about the data:

\begin{enumerate}
  \item \textbf{Ratings represent user preferences}: We assume that user ratings reflect their true preferences 
  (no gaming or manipulation).
  
  \item \textbf{Collaborative signal exists}: We assume that similar users rate similar movies similarly. 
  This underpins the collaborative filtering approach.
  
  \item \textbf{Data is clean}: The dataset is curated and does not require extensive cleaning. In production, 
  data validation and quality checks would be essential.
  
  \item \textbf{Static snapshot}: The dataset is treated as a static snapshot. In production, the training 
  data would be continuously refreshed.
  
  \item \textbf{No side information}: We use only ratings. User demographics, movie metadata, and tags are 
  not utilized (though they could improve predictions).
\end{enumerate}

\subsection{Sparsity and Its Implications}

The sparsity of the rating matrix is the central challenge in recommendation systems:

\[
  \text{Sparsity} = 1 - \frac{|\mathcal{R}|}{|\mathcal{U}| \times |\mathcal{I}|} = 1 - \frac{100,000}{943 \times 1,682} \approx 93.7\%
\]

Sparsity implies that:
\begin{itemize}
  \item Most user-item pairs have no observed rating.
  \item The rating matrix cannot be inverted or directly analyzed using standard matrix methods.
  \item Extrapolation from limited information (latent factors) is necessary.
  \item Cold-start (new users or items with few ratings) is a persistent problem.
\end{itemize}

Collaborative filtering through matrix factorization addresses sparsity by learning low-rank latent 
representations (Section~\ref{sec:methodology}).
